\subsection*{最初に必要な用語}
\addcontentsline{toc}{subsection}{最初に必要な用語}
\par\smallskip\noindent\refstepcounter{table}\label{tab:ch11-01}表 \thetable: 最初に必要な用語の整理\par\smallskip
表 \thetable は、最初に必要な用語の整理を比較・確認しやすい形で整理したものである。\par\smallskip
\begin{longtable}{>{\raggedright\arraybackslash}p{0.27\textwidth}>{\raggedright\arraybackslash}p{0.68\textwidth}}
\toprule
用語 & 初学者向けの説明 \\
\midrule
ソフトウェア工学 & ソフトウェアを、経験だけでなく、体系的・定量的・再現可能な方法で開発、運用、保守する分野である。\\
モデル & 対象の重要な側面だけを取り出して表したもの。図、表、数式、文章、形式仕様などの形をとる。\\
モデリング & モデルを作り、読み、分析し、関係者と共有する活動である。\\
方法 & 開発を進めるための体系化された進め方である。何を作り、どの順序で確認するかを含む。\\
抽象化 & 細部をあえて省き、重要な点だけを残すこと。モデルの基本である。\\
利害関係者 & ソフトウェアの開発結果に関心を持つ人や組織。利用者、顧客、開発者、認証機関、運用担当などである。\\
構文 & 表現が文法として正しいかを決める規則である。\\
意味論 & 表現が何を意味するかを決める規則である。\\
語用論 & 表現が文脈の中でどのように使われ、どう伝わるかに関する考え方である。\\
検証 & 作ったものが仕様や要求に合っているかを確認すること。\\
\bottomrule
\end{longtable}
